\section{Related Work}\label{related-work}

\subsection{Ontology Engineering and Knowledge Graph Construction}\label{ontology-engineering-and-knowledge-graph-construction}

\paragraph{Traditional Ontology Construction: Expressiveness vs.~Deployment Cost}\label{traditional-ontology-construction-expressiveness-vs.-deployment-cost}

The Semantic Web vision articulated by Berners-Lee, Hendler, and Lassila\cite{ref1} established the foundational goal of machine-readable, interlinked knowledge representations on the web. This vision was operationalized through the OWL 2 Web Ontology Language, whose formal semantics rooted in description logics provide expressiveness ranging from the lightweight OWL 2 EL profile to the fully decidable OWL 2 DL\cite{ref2}. Standard toolchains including the Prot\textquotesingle eg\textquotesingle e ontology editor\cite{ref3}, Apache Jena with its TDB triple store\cite{ref4}, and Virtuoso Universal Server\cite{ref5} enable ontology authoring, persistence, and SPARQL query evaluation at considerable scale. Reasoners such as HermiT and Pellet support automated classification and consistency checking, albeit with computational costs that scale poorly for large or rapidly evolving knowledge bases\cite{ref6}.

However, the deployment cost of full OWL-based pipelines remains prohibitive for many operational contexts. Triple stores require dedicated infrastructure; OWL reasoners impose batch-oriented computation ill-suited to streaming or event-driven workloads; and the expertise required for ontology engineering creates a bottleneck distinct from application development\cite{ref7}. These factors collectively confine heavyweight Semantic Web technologies to domains where upfront investment in formal knowledge representation is justified by long-term reasoning requirements, leaving a gap for lightweight, operational ontologies that can evolve at the speed of software engineering.

\paragraph{Lightweight Ontology Approaches: Accessibility Gains}\label{lightweight-ontology-approaches-accessibility-gains}

In response to these deployment barriers, the community has developed a spectrum of lightweight approaches that trade full logical expressiveness for accessibility and operational simplicity. Schema.org, launched in 2011 as a collaboration among major search engines, provides a pragmatic, property-centric vocabulary that now appears on over 40\% of web pages\cite{ref8}. Its design intentionally prioritizes decentralized adoption over semantic perfection, accepting property inheritance irregularities in exchange for rapid vocabulary evolution\cite{ref9}. JSON-LD\cite{ref10} extends conventional JSON with namespace support and URI resolution, allowing RDF-compatible serialization without requiring consumers to adopt full Linked Data processing pipelines. The Shapes Constraint Language (SHACL)\cite{ref11} provides data validation without inference, enabling closed-world constraint checking over RDF graphs that complements rather than replaces OWL's open-world semantics\cite{ref12}. Recent work has extended SHACL to support validation under updates\cite{ref13}, explainable validation through retrieval-augmented generation\cite{ref14}, and schema reasoning tasks\cite{ref12}, broadening its applicability beyond static integrity checking.

\textbf{Comparison with shape constraint languages.} SHACL and its alternative ShEx (Shape Expressions)\cite{ref48} represent the state of the art in RDF constraint validation and provide natural points of comparison for ADL Lite's \texttt{Comparator}-based precondition system. Table 2 contrasts the two families of approaches across eight salient dimensions. SHACL targets full RDF graphs with 50+ built-in constraint types (including path-based cardinality, value range, pattern matching, and logical combinations), validated through batch SPARQL queries or dedicated engines at \(O(|G|)\) cost per shape, where \(|G|\) is the graph size\cite{ref11}. ShEx provides a compact grammar-based syntax for describing RDF node shapes and their neighbor constraints, with validation through batch processing or API calls at comparable asymptotic cost\cite{ref48}. Both assume RDF as the underlying data model and require graph-processing infrastructure.

ADL Lite's \texttt{Comparator} enum---supporting scalar comparisons (\texttt{EQ}, \texttt{GT}, \texttt{LT}, \texttt{IN}, \texttt{NOT\_IN}, \texttt{REGEX}) over L1 YAML front matter fields---is deliberately simpler. It trades the full expressiveness of path-based graph constraints for \(O(1)\) per-rule validation cost at parse time through Pydantic, zero external dependencies, and operation within plain-text Markdown documents rather than RDF graphs. The design philosophy is complementary rather than competitive: ADL Lite uses the \texttt{Comparator} for lightweight scalar preconditions within document-native ontologies, while Phase 3 integration will layer SHACL for graph-level structural constraints over exported L3 relation triples. The \texttt{Comparator} handles \emph{scalar field validation} (e.g., ensuring \texttt{confidence\ \textgreater{}\ 0.7} before a \texttt{VALIDATE} action); SHACL handles \emph{graph structure validation} (e.g., ensuring every \texttt{Concept} has exactly one \texttt{status} property and at least one \texttt{actor} association).

\begin{longtable}[]{@{}>{\raggedright\arraybackslash}p{(\linewidth - 6\tabcolsep) * \real{0.2500}}
  >{\raggedright\arraybackslash}p{(\linewidth - 6\tabcolsep) * \real{0.2500}}
  >{\raggedright\arraybackslash}p{(\linewidth - 6\tabcolsep) * \real{0.2500}}
  >{\raggedright\arraybackslash}p{(\linewidth - 6\tabcolsep) * \real{0.2500}}@{}}
\toprule
Dimension
& ADL Lite \texttt{Comparator}
& SHACL
& ShEx
\\
\midrule
\endhead
\bottomrule
\endlastfoot
\textbf{Target} & L1 YAML front matter fields & RDF graphs & RDF nodes \\
\textbf{Constraint types} & Scalar comparison (\texttt{EQ}, \texttt{GT}, \texttt{LT}, \texttt{IN}, \texttt{NOT\_IN}, \texttt{REGEX}) & 50+ built-in constraints (\texttt{sh:minCount}, \texttt{sh:pattern}, \texttt{sh:node}, etc.) & Shape expressions (triple constraints, cardinality, recursion) \\
\textbf{Graph constraints} & No & Yes (path-based, cardinality, logical combinations) & Yes (neighbor constraints, reference shapes) \\
\textbf{Closed-world assumption} & Yes (by design) & Yes (\texttt{sh:closed}) & Yes (closed shapes) \\
\textbf{Validation timing} & Parse-time (Pydantic) & Batch / SPARQL endpoint & Batch / API \\
\textbf{Error reporting} & Pydantic field-level exceptions & \texttt{sh:ValidationReport} (RDF) & Typing errors (structural) \\
\textbf{Runtime cost} & \(O(1)\) per rule & \(O(|G|)\) per shape & \(O(|G|)\) per shape \\
\textbf{Extensibility} & Edit YAML registry & Write new SPARQL constraints (\texttt{sh:sparql}) & Write new shapes \\
\textbf{Dependencies} & Zero (Pydantic built-in) & RDF library + SHACL engine & RDF library + ShEx processor \\
\end{longtable}

These lightweight technologies collectively lower the barrier to structured data publishing but remain tethered to the RDF data model. They require triple stores or graph databases for non-trivial workloads and assume that data are already available in RDF-compatible formats---assumptions that exclude the vast corpus of domain knowledge encoded in Markdown documentation, plain-text notes, and collaborative wikis.

\paragraph{Document-Native Semantic Layers: Markdown-Based Knowledge Representation}\label{document-native-semantic-layers-markdown-based-knowledge-representation}

The past decade has witnessed significant growth in Markdown-based personal and collaborative knowledge management tools that implicitly encode graph structure through bi-directional linking. Foam\cite{ref15}, Obsidian\cite{ref16}, and Dendron\cite{ref17} extend plain-text Markdown with wiki-style links, backlinks, tag hierarchies, and graph visualization, enabling knowledge workers to construct personal knowledge graphs without leaving their editing environment. Dendron, in particular, introduces hierarchical note schemas with refactoring capabilities that preserve link integrity during structural changes\cite{ref18}, anticipating the need for schema-level governance in document-native knowledge bases.

These tools demonstrate that meaningful semantic structure can emerge from authoring conventions rather than formal ontological commitments. However, they lack machine-checkable coordination primitives: no constraint language prevents schema drift, no access control governs scope visibility, and no consensus mechanism mediates conflicting contributions from multiple agents. ADL Lite builds directly on this document-native paradigm while addressing each of these gaps.

\subsection{Palantir Foundry Data Engine}\label{palantir-foundry-data-engine}

\paragraph{FDE Ontology Layer: Nouns, Verbs, and Operational Semantics}\label{fde-ontology-layer-nouns-verbs-and-operational-semantics}

Palantir Foundry represents the most prominent industrial embodiment of an operational ontology platform. Its Data Engine exposes four core ontological primitives: Object Type (classes of entities), Property (attributes), Link (typed relationships), and Action (kinetic operations that modify the world)\cite{ref19}. This architecture unifies what Palantir terms ``semantic elements'' (nouns: objects, properties, links) with ``kinetic elements'' (verbs: actions, functions, dynamic security), enabling AI agents to both query and act upon the ontology through governed interfaces\cite{ref20}. The Ontology Metadata Service (OMS) defines all ontological entities, while Object Storage V2 supports tens of billions of objects per type with incremental indexing\cite{ref21}.

\paragraph{Digital Twin Model: Platform-Managed Object Representations}\label{digital-twin-model-platform-managed-object-representations}

Palantir's operational power derives from treating the ontology as a digital twin layer mediating between raw data and operational decisions\cite{ref21}. A dataset of GPS coordinates becomes a fleet tracking system when the ontology defines Vehicle objects, Route links, and PositionUpdate actions; a dataset of financial transactions becomes a risk monitoring system when the ontology defines Account objects, TransactionFlow links, and FraudAlert actions\cite{ref22}. Palantir's AIP (Artificial Intelligence Platform) operates within this ontology, with agents accessing objects, triggering actions, and reading relationships through governed interfaces rather than querying raw data\cite{ref21}.

\paragraph{Limitations: Enterprise Scale, Expert Authoring, and LLM-Agent Accessibility}\label{limitations-enterprise-scale-expert-authoring-and-llm-agent-accessibility}

Despite its sophisticated design, Palantir Foundry exhibits three limitations that motivate complementary approaches. First, deployment is enterprise-scale: the platform targets customers averaging millions of dollars in annual spend, with infrastructure requirements far exceeding those of Git-backed documentation or small-team knowledge bases\cite{ref19}. Second, ontology authoring requires specialized expertise and Palantir-specific tooling; Object Types, Link Types, and Action Types are configured through proprietary UIs rather than plain-text formats accessible to LLM agents or generalist developers\cite{ref20}. Third, the platform is not designed for LLM-native authorship: agents consume and act upon the ontology but do not author or evolve ontological definitions through conversational interfaces. ADL Lite addresses each limitation by offering pip-installable deployment, Markdown-native authoring with YAML-embedded semantics, and explicit support for multi-agent concept consensus through append-only event chains.

\subsection{Multi-Agent Consensus and Auditability}\label{multi-agent-consensus-and-auditability}

\paragraph{Consensus Mechanisms in Multi-Agent Systems}\label{consensus-mechanisms-in-multi-agent-systems}

Multi-agent systems require coordination mechanisms to reconcile divergent beliefs or proposals among autonomous participants. Classical approaches span voting protocols, reputation-based trust aggregation, and Byzantine fault-tolerant (BFT) consensus derived from distributed systems\cite{ref23}. Committee-based BFT protocols such as Practical Byzantine Fault Tolerance (PBFT) and its successors achieve agreement among partially trusted participants with logarithmic communication complexity, but assume a fixed membership model and synchrony bounds that limit applicability to open, dynamic agent collectives\cite{ref24}. Lightweight variants---including Proof-of-Contribution\cite{ref25}, utility-based Raft\cite{ref26}, and Comprehensive BFT (CBFT)\cite{ref24}---reduce consensus overhead by electing committees, partitioning storage, or adapting consensus granularity to network conditions. ADL Lite adapts these insights to a document-native setting: rather than achieving full BFT agreement over distributed state, it records individual agent proposals in an append-only chain and derives concept status through domain-specific quorum rules encoded in declarative YAML.

\paragraph{Append-Only Audit Logs for Agent Traceability}\label{append-only-audit-logs-for-agent-traceability}

The foundational work of Crosby and Wallach on tamper-evident logging introduced hash-chain data structures that provide \(O(1)\) append cost and \(O(j-i)\) verification for subsequences, detecting modification, deletion, and reordering attacks against untrusted log servers\cite{ref27}. These structures underpin Certificate Transparency (RFC 6962)\cite{ref28}, blockchain systems, and emerging agent audit standards. The recently proposed IETF Agent Audit Trail (AAT) draft specifies SHA-256 hash-chained JSON records with mandatory fields for agent identity, action classification, and trust level reporting, explicitly addressing EU AI Act Article 12 requirements for automatic event recording\cite{ref29}. Complementary proposals from industry---including append-only governance audit ledgers with cross-party verifiability\cite{ref30} and semantic memory ledgers such as MentisDB\cite{ref31}---extend tamper-evident logging from system events to agent cognition, recording not merely actions but decisions, corrections, constraints, and handoffs as durable cognitive checkpoints. ADL Lite's EventChain applies this architectural pattern to ontological evolution: each concept exists as a cryptographically linked sequence of status-changing events, rendering the entire lifecycle of a definition---from proposal through challenge to resolution---forensically reconstructible.

\paragraph{AML Pattern Detection as Multi-Agent Evaluation Domain}\label{aml-pattern-detection-as-multi-agent-evaluation-domain}

Anti-money laundering (AML) pattern detection provides a demanding evaluation domain for multi-agent ontological consensus due to the complexity of financial crime typologies, the heterogeneity of data sources, and the severe consequences of conceptual error. Financial institutions increasingly employ knowledge graphs---connecting people, companies, accounts, and transactions---to uncover hidden relationships, identify shell company structures, and trace beneficial ownership across jurisdictional boundaries\cite{ref32}. Semantic rule-based approaches model AML domain knowledge as OWL ontologies with Jena inference engines, executing threshold-based rules over transaction data to flag suspicious patterns\cite{ref33}. These systems face challenges of data quality, scalability, and interpretation complexity that directly motivate machine-readable, auditable concept definitions\cite{ref32}. ADL Lite's pilot evaluation targets this domain not as a replacement for production AML platforms, but as a proving ground for ontological coordination: the precision required for financial crime concepts imposes strict correctness requirements that expose weaknesses in informal or mutable knowledge representation approaches.

\paragraph{Provenance and Verifiable Publishing Standards}\label{provenance-and-verifiable-publishing-standards}

The landscape of Semantic Web provenance and verifiable publishing encompasses several standards that address complementary aspects of the problem space that ADL Lite targets: recording \emph{who} said \emph{what}, \emph{when}, and with \emph{what} evidentiary basis. This subsection surveys four standards---W3C PROV-O, nanopublications, Linked Data Proofs (now W3C Verifiable Credential Data Integrity), and RDF-star---that collectively define the state of the art in provenance-aware knowledge representation, and positions ADL Lite's EventChain model with respect to each.

\textbf{W3C PROV-O.} The W3C PROV Ontology \cite{ref34} provides a standard model for provenance interchange on the Web, expressing the PROV Data Model \cite{ref35} in OWL 2. PROV-O defines core concepts including \texttt{Entity} (the subject of provenance), \texttt{Activity} (a process that generates or transforms entities), and \texttt{Agent} (an entity responsible for an activity), along with properties such as \texttt{wasGeneratedBy}, \texttt{wasAssociatedWith}, \texttt{wasAttributedTo}, and \texttt{wasInformedBy} that establish causal and responsibility relationships among provenance elements. PROV-O has been adopted across scientific domains including bioinformatics \cite{ref36}, materials science \cite{ref37}, and workflow provenance systems \cite{ref38}, and has been extended by specialized profiles such as ProvONE for scientific workflow provenance \cite{ref38} and D-PROV for structured data transformations \cite{ref39}.

ADL Lite's EventChain can be directly mapped to PROV-O: each Event in an EventChain corresponds to a \texttt{prov:Activity}, the concept being modified corresponds to a \texttt{prov:Entity}, and the \texttt{actor} field in each event maps to \texttt{prov:wasAssociatedWith}. The \texttt{previous\_event\_id} cryptographic linkage provides a \texttt{prov:wasInformedBy} chain connecting successive activities. However, PROV-O alone does not provide cryptographic integrity guarantees---it is a data model for describing provenance, not a mechanism for verifying it. An RDF graph serialized in PROV-O can be modified, re-serialized, and redistributed without detection unless additional signing or hashing infrastructure is layered on top. ADL Lite complements PROV-O by extending its descriptive provenance model with built-in SHA-256 hash chaining: the EventChain is simultaneously a PROV-O-compatible provenance trace and a tamper-evident data structure whose integrity can be verified independently of any external trust infrastructure. Future Phase 3 work includes a PROV-O export serializer that will materialize ADL Lite EventChains as fully standard-compliant RDF provenance graphs.

\textbf{Nanopublications and Trusty URIs.} The nanopublication framework \cite{ref40} provides a standardized way to publish small, granular scientific claims with cryptographically signed provenance. Each nanopublication consists of three named graphs \cite{ref41}: an assertion graph (the claim being made), a provenance graph (metadata about the claim's origin), and a publication-info graph (metadata about the nanopublication itself). The Trusty URI mechanism \cite{ref42} provides content-addressed identifiers via Base64-encoded SHA-256 hashes, ensuring that any modification to a nanopublication's content changes its URI and thus invalidates the published artifact. This achieves decentralized verifiable publishing---a goal shared with ADL Lite's append-only EventChain---where published claims are immutably bound to their provenance and can be independently verified without reliance on a central authority.

Key differences between the two approaches reveal complementary rather than competing design goals. First, nanopublications focus on \emph{atomic claims} with single assertions (e.g., ``gene X is associated with disease Y''), while ADL Lite models full \emph{concept lifecycles} as multi-event chains spanning registration, validation, deprecation, forking, and archival. Second, nanopublications use RSA digital signatures for author authentication, binding each assertion to a verifiable cryptographic identity; ADL Lite Phase 1 currently relies on string identifiers for actors, with LD-Proof integration planned for Phase 3 to achieve equivalent signature-based authentication. Third, nanopublications require RDF/SPARQL infrastructure for publication, querying, and verification, while ADL Lite operates over Markdown files in standard Git repositories without a triple store. Nanopublications thus target the publication of individual scientific assertions to decentralized servers, while ADL Lite targets the governance of evolving concept definitions within multi-agent document-native workflows.

\textbf{Linked Data Proofs and RDF Dataset Canonicalization.} The W3C Verifiable Credential Data Integrity specification \cite{ref43}---formerly known as Linked Data Proofs---describes mechanisms for ensuring the authenticity and integrity of constrained digital documents using cryptography, especially through digital signatures over canonicalized RDF datasets. The specification defines a vocabulary for proof types, verification methods, and canonicalization algorithms, and is accompanied by concrete cryptographic suites including EdDSA (Ed25519Signature2020) and ECDSA variants \cite{ref44}. Underpinning this signing capability is the RDF Dataset Canonicalization algorithm (URDNA2015) \cite{ref45}, which produces a deterministic, byte-for-byte consistent serialization of an RDF dataset regardless of the original ordering of triples or blank node labels---a prerequisite for signature verification over RDF data.

This technology provides the authenticated provenance that ADL Lite currently lacks: the ability to bind actor identities to events through cryptographic signatures rather than string identifiers. An LD-Proof over an ADL Lite event would enable any verifier to confirm that a specific actor (identified by a Decentralized Identifier, or DID) authored a specific event, without requiring trust in the storage medium or the party presenting the chain. The reviewer specifically recommended this as a path forward for authenticated event logging in ADL Lite. The current status of LD-Proofs is a W3C Recommendation (May 2025), with implementations including \texttt{vc-js} and \texttt{jsonld-signatures} that support multiple cryptographic suites. ADL Lite's future Phase 3 includes LD-Proof integration for authenticated events, targeting \texttt{Ed25519Signature2020} as the initial suite given its balance of security and computational efficiency.

\textbf{RDF-star and Named Graphs.} RDF-star \cite{ref46}---now incorporated into the RDF 1.2 specification \cite{ref47}---provides statement-level annotation by enabling triples to function as subjects or objects of other triples (so-called ``quoted'' or ``embedded'' triples), without the syntactic verbosity of explicit reification nodes or the artificial graph proliferation of single-triple Named Graphs. The annotation syntax \texttt{\{\textbar{}\ :source\ :bob\ \textbar{}\}} in Turtle-star allows concise expression of metadata about individual statements, such as provenance, certainty, or temporal validity, directly adjacent to the statement being annotated. Named Graphs \cite{ref41}, standardized as part of SPARQL 1.1, allow bundling triples into named containers with their own provenance metadata, supporting graph-level rather than statement-level provenance tracking.

Both technologies address the ``who said what when'' problem that ADL Lite's EventChain solves through a different architectural path. RDF-star enables statement-level annotations within a single graph; Named Graphs enable graph-level provenance partitioning; ADL Lite's EventChain enables event-level provenance sequencing with cryptographic integrity guarantees. ADL Lite's L3 relation blocks---which encode \texttt{source}-\texttt{relation}-\texttt{target} triples within Markdown fenced code blocks---can be viewed as a Markdown-native form of RDF triples, with each EventChain entry providing statement-level provenance analogous to an RDF-star annotation. The key distinction is that RDF-star annotations are syntactic constructs within an RDF graph (with no built-in integrity protection), while ADL Lite EventChain entries are cryptographically linked events in an append-only sequence (with built-in tamper detection but no native RDF serialization in Phase 1). Phase 3 work will include an RDF-star export serializer that materializes ADL Lite relation blocks as quoted triples with EventChain-derived annotations, bridging these two provenance models.

\paragraph{Event-Centric Ontologies}\label{event-centric-ontologies}

The ontologies surveyed in the preceding subsections vary in their treatment of events: some model events as first-class entities central to their philosophical commitment, while others treat events as secondary constructs or ignore them entirely. This subsection examines three event-centric approaches---CIDOC CRM, SEM/LODE, and RDF Stream Processing---that share ADL Lite's fundamental premise that \emph{change in the world is best understood through events} rather than through direct property mutations of persistent objects, and positions ADL Lite's EventChain model within this intellectual lineage.

\textbf{CIDOC CRM (Conceptual Reference Model).} The CIDOC Conceptual Reference Model\cite{ref49} is an ISO 21127 standard for cultural heritage information integration, developed by the International Committee for Documentation of the International Council of Museums (ICOM-CIDOC). CIDOC CRM is fundamentally event-centric: it models reality through the \texttt{E5\ Event} class that brings together \texttt{E39\ Actor} participants with \texttt{E18\ Physical\ Thing} objects in spatiotemporal contexts. The ontology's core philosophical commitment---expressed through its formal axioms and extensive documentation---is that change in the world is understood through events, not through direct property mutations of persistent objects. An actor painting a canvas, a museum acquiring an artifact, or an archaeological excavation uncovering a fragment are all modeled as events with their own identity, temporal extent, spatial location, and participant roles, rather than as implicit state changes to the participating entities.

This ontological stance directly parallels ADL Lite's design: ADL Lite models concept lifecycle changes (registration, validation, deprecation, archival) as explicit events in an append-only chain rather than as mutable field updates to a concept record. A concept's \texttt{status} is not a stored property but a value \emph{derived} from the history of events affecting that concept, mirroring CIDOC CRM's treatment of an object's current custodian as derivable from the sequence of custody-transfer events. ADL Lite's \texttt{previous\_event\_id} cryptographic linkage serves an analogous function to CIDOC CRM's temporal and causal ordering axioms, ensuring that the event sequence is both immutable and verifiable.

The key difference lies in purpose and scope. CIDOC CRM is \emph{descriptive}: it documents what happened in the past (or what is believed to have happened), providing a comprehensive upper ontology for cultural heritage documentation with 81 classes and over 200 properties. ADL Lite is \emph{operational}: it enforces what \emph{can} happen through declarative preconditions (\texttt{PreconditionRule}), governing the future evolution of concept definitions rather than documenting their past. CIDOC CRM supports inference over historical event patterns; ADL Lite supports runtime validation of proposed action transitions. ADL Lite shares CIDOC CRM's philosophical commitment to event centrality but applies it to a different domain---operational ontology lifecycle governance rather than cultural heritage documentation---and with a different deployment model---plain-text Markdown files rather than RDF triple stores.

\textbf{SEM (Simple Event Model) and LODE (Linked Open Descriptions of Events).} While CIDOC CRM provides a comprehensive upper ontology, the Simple Event Model (SEM)\cite{ref50} and Linked Open Descriptions of Events (LODE)\cite{ref51} represent lightweight alternatives designed for publishing event data on the Semantic Web. SEM provides a minimal vocabulary comprising \texttt{sem:Event}, \texttt{sem:Actor}, \texttt{sem:Place}, and \texttt{sem:Time} for describing events without committing to a heavy upper ontology. Its design principle is epistemic modesty: SEM makes minimal ontological commitments, allowing publishers to describe what they know about an event without requiring complete provenance or exhaustive participant enumeration. LODE extends SEM with Dublin Core compatibility (\texttt{dc:title}, \texttt{dc:description}, \texttt{dc:date}) and SKOS-based event type hierarchies, enabling integration with existing Linked Data publishing workflows\cite{ref51}.

Both SEM and LODE are \emph{descriptive} ontologies for event publishing: they enable data providers to expose event descriptions in RDF for consumption, linking, and querying across the Web of Data. They do not provide constraint mechanisms for governing which events can occur, nor do they offer integrity guarantees for event sequences. ADL Lite complements these descriptive frameworks with an \emph{operational} layer: the EventChain is not merely a published description of events but a tamper-evident, precondition-governed mechanism for controlling concept lifecycle transitions. An SEM publisher might describe that a concept was validated by agent X at time T; ADL Lite's EventChain \emph{enforces} that the validation can occur only if the concept is in \texttt{PROPOSED} status, the agent's scope permits validation actions, and all preconditions evaluate to true---then cryptographically binds the event into an immutable sequence. SEM and LODE enable event data interoperability; ADL Lite enables event-driven concept governance with cryptographic integrity.

\textbf{RDF Stream Processing (RSP).} The W3C RDF Stream Processing Community Group defined models and query languages for continuous processing over timestamped RDF triples arriving in streams\cite{ref52}. Systems such as C-SPARQL\cite{ref53} and CQELS\cite{ref54} extend SPARQL with windowing operators (\texttt{RANGE}, \texttt{STEP}) and continuous query semantics, enabling real-time pattern detection over streaming RDF data. In RSP systems, events are modeled as timestamped triples (or triple graphs) arriving in potentially unbounded streams; queries register continuous patterns (e.g., ``report all transactions exceeding $10,000$ within a 5-minute window'') that are evaluated incrementally as new data arrives.

ADL Lite's EventChain can be viewed as a \emph{materialized} RDF stream: each event appends triples (encoded as YAML front matter fields and L3 relation blocks) to a persistent, cryptographically linked chain, enabling retrospective provenance queries rather than real-time pattern detection. Where C-SPARQL processes streaming triples in temporal windows for immediate alerting, ADL Lite processes historical events in an append-only chain for forensic reconstruction and status derivation. This distinction reflects a fundamental difference in target workload: RSP systems optimize for low-latency detection over high-velocity streams, while ADL Lite optimizes for integrity and auditability over low-velocity, high-stakes concept lifecycles. Nevertheless, the RSP formalization of events as timestamped RDF triples provides a conceptual bridge to future Phase 3 work: an RSP-compatible export of ADL Lite EventChains would enable continuous query patterns over concept lifecycle events (e.g., ``alert when any concept receives three \texttt{CHALLENGE} events within 24 hours''), extending retrospective audit to proactive monitoring.

\paragraph{Conflict-Free Replicated Data Types for Collaborative Ontology Authoring}\label{conflict-free-replicated-data-types-for-collaborative-ontology-authoring}

Multi-agent ontology authoring introduces a fundamental concurrency challenge: when multiple agents modify concept definitions simultaneously---whether through concurrent commits in separate Git branches or through simultaneous edits in a real-time collaborative session---their changes must be reconciled without losing information or violating consistency. Conflict-Free Replicated Data Types (CRDTs)\cite{ref55} provide a theoretical and practical foundation for this reconciliation, ensuring that concurrent updates converge to a consistent state without requiring centralized coordination or locking. This subsection surveys CRDT literature relevant to ADL Lite's multi-agent authoring model and identifies opportunities for future integration.

\textbf{Automerge.} Automerge\cite{ref56} is a JSON CRDT library that provides automatic merge semantics for collaborative document editing. Automerge uses operation-based CRDTs with a novel ``multi-value register'' semantics that preserves all concurrent edits rather than choosing a single winner. When two agents concurrently modify the same field---for example, one agent updates a concept's \texttt{description} while another updates its \texttt{scope}---Automerge merges both changes automatically, preserving the intent of both editors. In cases of true conflict (both agents modify the \emph{same} field concurrently), Automerge preserves both values as a multi-value register rather than silently applying last-write-wins (LWW), enabling downstream resolution by application logic or human review.

This behavior is directly relevant to ADL Lite's fork model. When two agents concurrently edit the same concept---for instance, both proposing alternative \texttt{description} values for a concept in \texttt{CHALLENGED} status---ADL Lite's current Phase 1 implementation applies last-write-wins to the event chain, potentially discarding one agent's contribution. Automerge would preserve both versions as concurrent events, allowing the multi-agent quorum mechanism to evaluate both proposals rather than arbitrarily rejecting one. Adopting CRDT semantics for event chain merging is a promising future direction: ADL Lite EventChain entries could be modeled as operations in a sequential CRDT, with the \texttt{previous\_event\_id} hash chain providing causal ordering and the CRDT layer providing automatic convergence for concurrent appends. This hybrid approach---cryptographic integrity from hash chaining plus automatic merge semantics from CRDTs---would eliminate the LWW limitation without sacrificing tamper evidence.

\textbf{Yjs.} Yjs\cite{ref57} is a widely deployed CRDT framework for real-time collaborative editing, optimized for shared text documents. Yjs provides efficient delta encoding, incremental update propagation, and awareness protocols (cursor positions, user presence) for concurrent editing of textual content. Its underlying YATA algorithm\cite{ref58} for list CRDTs achieves sub-linear space overhead for concurrent insertions and deletions, making it suitable for large documents with many concurrent editors.

While Yjs targets text documents, its underlying CRDT principles are applicable to event sequence merging in ADL Lite. An ADL Lite EventChain is, at its core, an ordered sequence of event entries---structurally analogous to a list of text operations. The YATA algorithm's list CRDT semantics could be adapted to merge concurrent event sequences: when two agents append events to the same concept's chain in different branches, the YATA merge would produce a convergent ordering that preserves all events and their causal relationships. The primary adaptation required would be extending YATA's content-based ordering with ADL Lite's hash-based causal linking (\texttt{previous\_event\_id}), ensuring that cryptographic integrity and CRDT convergence operate in concert. This integration remains future work but demonstrates that ADL Lite's document-native, Git-backed architecture is compatible with state-of-the-art collaborative editing technology.

\textbf{CRDT-Ontology Intersection.} Recent research has explored applying CRDTs to knowledge graph editing, enabling concurrent multi-agent ontology evolution without centralized coordination. Chesa et al.\cite{ref59} demonstrated a CRDT-based approach to collaborative RDF editing where concurrent property updates converge automatically while preserving graph consistency invariants. Parallel work on distributed ontology editing platforms\cite{ref60} has explored CRDTs for merging concurrent class definition changes, property range modifications, and instance assertions in OWL ontologies. These approaches share ADL Lite's motivation---enabling multiple agents to evolve an ontology concurrently---but target full RDF/OWL graph editing rather than ADL Lite's lightweight Markdown-native event chains.

ADL Lite's Git-native design already provides branch-based concurrent authoring: agents work in isolated branches, propose changes through event chains, and merge through pull requests. What CRDTs would add is \emph{automatic merge semantics for event chains}: rather than requiring manual conflict resolution when branches diverge, a CRDT-backed EventChain would converge automatically while preserving the cryptographic integrity guarantees of the hash chain. The envisioned integration path is a hybrid model where (1) each branch maintains a local CRDT-backed event sequence, (2) cross-branch merges use CRDT convergence to produce a unified ordering, and (3) the resulting merged sequence is cryptographically hashed to produce a tamper-evident canonical chain. This would extend ADL Lite's current Git-based concurrency model with real-time collaborative capabilities while maintaining the forensic traceability that is central to its design.

\subsection{Positioning}\label{positioning}

\paragraph{Comparison Across Seven Dimensions}\label{comparison-across-seven-dimensions}

\textbf{Table 1. Positioning of ADL Lite} against three reference approaches---Palantir FDE, OWL/SPARQL pipelines, and Markdown-only knowledge bases---across seven dimensions salient to operational ontology deployment.

\begin{longtable}[]{@{}>{\raggedright\arraybackslash}p{(\linewidth - 8\tabcolsep) * \real{0.1667}}
  >{\centering\arraybackslash}p{(\linewidth - 8\tabcolsep) * \real{0.2083}}
  >{\centering\arraybackslash}p{(\linewidth - 8\tabcolsep) * \real{0.2083}}
  >{\centering\arraybackslash}p{(\linewidth - 8\tabcolsep) * \real{0.2083}}
  >{\centering\arraybackslash}p{(\linewidth - 8\tabcolsep) * \real{0.2083}}@{}}
\toprule
Dimension
& Palantir FDE
& OWL/SPARQL
& Markdown-only
& ADL Lite
\\
\midrule
\endhead
\bottomrule
\endlastfoot
\textbf{Deployment Cost} & Enterprise (\(M\)+) & High (triple store + reasoner) & Low (Git repo) & Very Low (\texttt{pip\ install} + Git) \\
\textbf{LLM Authorship} & Not designed & Requires ontology expertise & Informal (free-text) & Native (YAML + Markdown) \\
\textbf{Integrity Guarantees} & Platform-managed & Reasoner-based consistency & None & SHA-256 EventChain (append-only) \\
\textbf{Action Preconditions} & TypeScript functions & OWL/SWRL rules & None & Declarative YAML + SSA validation \\
\textbf{Consensus Mechanism} & Centralized (admin UI) & N/A (single authority) & N/A (no coordination) & Multi-agent quorum + event history \\
\textbf{Scalability} & Billions of objects & Millions of triples & Thousands of notes & Pilot: 20 concepts; target: \(10^4\) \\
\textbf{Standards Alignment} & Proprietary (Palantir OMS) & Full W3C stack (OWL, RDF, SPARQL, SHACL) & None & Partial (PROV-O mappable, SHACL-integrable, no native RDF serialization in Phase 1) \\
\end{longtable}

\paragraph{ADL Lite: RDF-Like Triples in Markdown, Machine-Checkable Coordination Without the Triple Store}\label{adl-lite-rdf-like-triples-in-markdown-machine-checkable-coordination-without-the-triple-store}

Palantir FDE delivers industrial-strength operational semantics through a tightly integrated platform, but at deployment cost and accessibility barriers that exclude small teams, open-source projects, and LLM-native workflows. OWL/SPARQL pipelines offer unmatched formal expressiveness through description logic semantics and automated reasoning, yet require specialized expertise and infrastructure that confine them to research and enterprise contexts with dedicated ontology engineering functions. Markdown-based tools democratize knowledge capture through plain-text authoring and Git-based collaboration, but provide no machine-checkable guarantees against schema drift, conflicting definitions, or unauthorized modification.

ADL Lite occupies a deliberate middle ground among these alternatives. L3 relation blocks provide RDF-like triples editable in ordinary Markdown; L1 YAML scalars encode datatype properties with zero overhead beyond the YAML front matter already conventional in static site generators and documentation tools. SSA (Status-State-Action) validation, scope ACL, and the append-only EventChain supply machine-checkable coordination primitives---status transitions with preconditions, visibility boundaries, and tamper-evident history---without requiring a triple store, OWL reasoner, or platform-specific infrastructure. The comparison is not one of replacement but of complementarity: ADL Lite targets contexts where the semantic precision of OWL and the operational power of Palantir are desirable but the deployment investment is disproportionate to the available resources.

This positioning has direct implications for multi-agent systems. Palantir's AIP demonstrates that agents can operate effectively within an ontology, but presupposes a fully staffed, platform-managed ontology layer\cite{ref21}. OWL supports agent reasoning through description logic services but assumes centralized ontological authority. Markdown imposes no authority structure but offers no coordination primitives at all. ADL Lite enables a new mode: multiple LLM agents proposing, challenging, and refining concepts within a shared document-native ontology, with the EventChain providing the audit trail and consensus mechanism that makes such decentralized authorship accountable. Phase 1 pilots (\(n=20\) AML concepts) demonstrate mechanistic gains in lifecycle traceability, scope isolation, and L3-sensitive retrieval rather than claims of automated ontological inference. We argue that a lightweight, document-native semantic layer---formalized in future work through a centralized OntologyManager schema---can deliver multi-agent auditability at deployment cost closer to Git-backed notes than enterprise knowledge graphs.

Compared to PROV-O-based provenance logs, ADL Lite provides built-in cryptographic integrity verification rather than provenance description alone; PROV-O excels at modeling \emph{what happened}, while ADL Lite's EventChain additionally guarantees \emph{that the record of what happened has not been tampered with}. Compared to nanopublications, ADL Lite targets concept-level lifecycle governance (registration through validation to deprecation and archival) rather than atomic claim publishing; nanopublications excel at immutable assertion dissemination, while ADL Lite excels at governing the evolution of definitions over time. Compared to CIDOC CRM, ADL Lite shares the philosophical commitment to event-centric modeling but targets operational enforcement of lifecycle transitions rather than descriptive documentation of cultural heritage events. Compared to SEM/LODE, ADL Lite provides precondition-governed event validation and cryptographic integrity rather than lightweight event description for Linked Data publishing. Compared to Git-only pipelines, ADL Lite adds structured action preconditions (the \texttt{Comparator}-based \texttt{PreconditionRule} system), status derivation from event history (computing \texttt{status} and \texttt{confidence} from the chain rather than storing them as mutable fields), and a closed action registry that eliminates runtime code injection. The gap that ADL Lite fills is the intersection of four properties that no existing approach combines: (a) document-native authoring in plain Markdown, (b) built-in tamper detection through cryptographic hash chaining, (c) lifecycle state machines with declarative preconditions, and (d) pip-installable deployment requiring no enterprise infrastructure. PROV-O provides (a) only when combined with RDF authoring tools and lacks (b) natively; nanopublications provide (b) for published claims but lack (a), (c), and (d); CIDOC CRM provides event-centric modeling but requires heavyweight RDF infrastructure; Git provides (a) and (d) but lacks (b) and (c); Palantir provides (b), (c), and (d) only at enterprise scale with proprietary tooling. It is this unique combination---document-native, tamper-evident, precondition-governed, and frictionlessly deployable---that motivates ADL Lite as a complementary contribution to the provenance-aware knowledge representation landscape.
